% Generated draft: requires mathematical and source review.
\documentclass[11pt]{article}
\usepackage[margin=25mm]{geometry}
\usepackage{fontspec}
\setmainfont{Latin Modern Roman}
\usepackage{amsmath,amssymb,mathtools,unicode-math}
\setmathfont{Latin Modern Math}
\usepackage{enumitem,needspace,xurl,hyperref}
\setlength{\parindent}{0pt}
\setlength{\parskip}{5pt}
\setlength{\emergencystretch}{3em}
\newcommand{\N}{\mathbb{N}}
\newcommand{\Z}{\mathbb{Z}}
\newcommand{\Q}{\mathbb{Q}}
\newcommand{\R}{\mathbb{R}}
\newcommand{\C}{\mathbb{C}}
\newcommand{\F}{\mathbb{F}}
\newcommand{\A}{\mathbb{A}}
\newcommand{\PP}{\mathbb{P}}
\newcommand{\E}{\mathbb{E}}
\newcommand{\eps}{\varepsilon}
\newcommand{\dd}{\,\mathrm{d}}
\begin{document}
\section*{UnsolvedMath ID: 30001859}
\textbf{Source problem code:} catalogue entry
\subsection{Definitions and mathematical statement}
For a positive integer \(n\), let \(C_n\) denote a cyclic group of order \(n\), and set
\[
G=C_2\times C_2\times C_2\times C_6.
\]
The integral group ring \(\Z G\) consists of sums \(\sum_{g\in G}a_g g\), with \(a_g\in\Z\), with coefficientwise addition and multiplication determined by the multiplication in \(G\) and distributivity. Its unit group is
\[
\mathrm{U}(\Z G)=\{u\in\Z G:\text{ there exists }v\in\Z G\text{ with }uv=vu=1\}.
\]
\textbf{Source conjecture, with unresolved terminology.} There exists a unit \(u\in\mathrm{U}(\Z G)\) that is an ``exotic unit'' in the sense intended by the cited report.

The supplied record does not define ``exotic unit'' or specify the class of units against which exoticity is measured. It supplies no additional augmentation, order, equivalence, or generation conditions on the sought unit. Consequently, the displayed existential formulation records the intended conjecture but does not provide a fully specified mathematical predicate. The supplied literature assessment also warns that the extracted group context is incomplete; the group displayed above is retained exactly as supplied, pending review of that context.
\subsection{Short English statement}
The source conjectures that the integral group ring of \(C_2\times C_2\times C_2\times C_6\) contains an exotic unit. The meaning of ``exotic'' and possibly relevant surrounding group context are missing, so a precise existence statement requires human review of the cited report. The record labels the problem open; this status is supplied curation, not an independently verified conclusion.
\subsection{Sources}
\begin{itemize}
\item Local registry: \texttt{data/unsolvedmath/problems.json}, permanent UnsolvedMath ID \(30001859\), source problem code \texttt{catalogue entry}; collection: Oberwolfach Reports Open Problems. The record gives proposed year 2011 and supplies no named proposer.
\item Supplied source citation: \emph{Computational Group Theory (2012)}, Oberwolfach Reports, DOI \texttt{10.4171/owr/2011/37}, pp.\ 2113--2161. \url{https://doi.org/10.4171/owr/2011/37}. The supplied background describes these reports as documenting mathematical workshops and their open-problem discussions.
\item Supplied literature source, labeled ``OWR report 37'': \url{https://doi.org/10.4171/owr/2011/37}. Literature checked at: \texttt{2026-08-21}. This is supplied dated curation, not independently verified. Its assessment labels the problem open, warns that the extracted group context is incomplete, describes the exotic-unit assertion as a conjecture, and reports that no authoritative resolution for the named group was found. The background explicitly requests independent verification of time-sensitive status and cited claims.
\item Supplied statement-verification metadata labels the statement \texttt{corrected\_verified} and says that the official report confirms the conjecture, that an extraction label was removed, and that the group notation was restored. These are claims of the supplied record; no independent source inspection is asserted here.
\end{itemize}
\end{document}
